\subsection{补算：移动磁通面的两项变化}
固定曲面上的磁通变化由Faraday定律给出
$-\int_S\grad\times\E\cdot\dd\vct S=-\oint_C\E\cdot\dd\vct l$。
若边界在$\dd t$内移动$\vct v_f\dd t$，新增的窄条面积贡献为
$-\dd t\oint_C(\vct v_f\times\B)\cdot\dd\vct l$；符号由曲面法向和边界右手定向决定。两者相加正是式\eqref{eq:flux}。
例如均匀$\B=B\hat z$中面积扩大的圆环，外移速度径向，$\vct v_f\times\B$沿负切向，因此公式给正磁通增量，可用来检查负号。

把$\E+\vct v_f\times\B=\vct R$记为所选流速下的非理想电场。保持所有随流回路磁通的充分局部条件是$\grad\times\vct R=0$，在单连通区域可写成单值势梯度。$\vct R$非零本身并不够；它可能只有静电梯度。反过来，局地磁通输运速度存在也不自动证明全局磁连接关系没有改变，奇点、边界与三维场线积分都需检查。

\subsection{标量电子压力的旋度：等温与非正压的区别}
利用$\grad\times(f\grad g)=\grad f\times\grad g$，
\begin{equation}
 \grad\times{\grad p_e\over en}
 =\grad\left({1\over en}\right)\times\grad p_e
 =-{\grad n\times\grad p_e\over en^2}.
\end{equation}
若$p_e=p_e(n)$，两个梯度平行，旋度为零。例如等温电子$p_e=nT_{e0}$时
\begin{equation}
 {\grad p_e\over en}={T_{e0}\over e}\grad\ln n
 =\grad\left({T_{e0}\over e}\ln{n\over n_*}\right).
\end{equation}
于是无惯性、无电阻的电子冻结仍成立。若$p_e=nT_e(\vct x)$，则
$\grad n\times\grad p_e=n\grad n\times\grad T_e$，不再保证为零。这是非正压源项，不能与电子非对角压力的作用混称。

\subsection{电子广义磁场的旋度运算，一步也不省}
在正压、无阻力条件下，用
$(\vel_e\cdot\grad)\vel_e=\grad(u_e^2/2)-\vel_e\times\vct\omega_e$
把电子欧姆定律写成
\begin{equation}
 \E+\vel_e\times\B
 =-\grad H_e-{m_e\over e}\partial_t\vel_e
 -{m_e\over e}\grad(u_e^2/2)
 +{m_e\over e}\vel_e\times\vct\omega_e.
\end{equation}
对全式取负旋度，用$\partial_t\B=-\grad\times\E$，梯度项的旋度为零：
\begin{equation}
 \partial_t\B
 =\grad\times(\vel_e\times\B)
 +{m_e\over e}\partial_t\vct\omega_e
 -{m_e\over e}\grad\times(\vel_e\times\vct\omega_e).
\end{equation}
将时间导数合并到左边，就得到
$\partial_t(\B-m_e\vct\omega_e/e)=
\grad\times[\vel_e\times(\B-m_e\vct\omega_e/e)]$。
这条广义输运关系来自电子方程本身；将它进一步化为
$\B_e=\B-d_e^2\nabla^2\B$，才另外要求固定均匀离子、常密度及低频Ampère关系。

常密度时，
\begin{align}
 \vct\omega_e
 &=-{1\over\mu_0en}\grad\times(\grad\times\B)\\
 &=-{1\over\mu_0en}[\grad(\grad\cdot\B)-\nabla^2\B]
 ={1\over\mu_0en}\nabla^2\B.
\end{align}
变密度时还多出$\grad(1/n)\times(\grad\times\B)$，不能把$d_e$改成一个空间函数后仍照用常系数式。

\subsection{Poynting定理和碰撞热：补齐总能量账本}
Faraday方程点乘$\B/\mu_0$，Ampère--Maxwell方程点乘$\E$，相加并用
\begin{equation}
 \grad\cdot(\E\times\B)=\B\cdot(\grad\times\E)-\E\cdot(\grad\times\B)
\end{equation}
即可得到前面的Poynting定理。与式\eqref{eq:energy}对物种求和后，$\sum_s q_sn_s\vel_s=\J$使$\J\cdot\E$严格抵消。保守碰撞还要求
\begin{equation}
 \sum_s(\vel_s\cdot\vct R_s+H_s)=0.
\end{equation}
仅有$\vct R_i+\vct R_e=0$只保证动量守恒，尚未保证能量守恒。

以$\vct R_e=\alpha(\vel_i-\vel_e)$、$\vct R_i=-\vct R_e$为例，
\begin{equation}
 \vel_e\cdot\vct R_e+\vel_i\cdot\vct R_i
 =-\alpha|\vel_i-\vel_e|^2=-\eta J^2.
\end{equation}
若碰撞在系统内部发生，就必须有$H_i+H_e=\eta J^2$回到内能；具体如何分给两物种需要碰撞模型。只加电阻扩散却不回收热量的算例，不是孤立总能量守恒模型。

Hall项则始终满足$\J\cdot(\J\times\B)/(en)=0$。不过单种粒子的电场功$q_sn_s\vel_s\cdot\E$可以非零，能流也可以被Hall结构重新导向；“总Hall功为零”和“Hall过程不影响加热位置与能量分配”不是同一判断。
